Can a deterministic multitape Turing machine transpose a $d\times d$ matrix with $\lceil\log_2d\rceil$-bit entries in
time linear in the input length~$L$? Sixteen earlier papers of this series study the question without deciding it. This
paper is a synthesis of them, a diagnosis of where they stop, and two new theorems.

The synthesis states the main results of the series in one place, with their main hypotheses. The diagnosis is
this. For oblivious machines the question is now one about a single function, the best rate of a storage code on the
bounded-depth cone relation that is left of a linear-time schedule after a genie has revealed $\eps L$ bits: a
transposer needs rate $1-\eps$ for every~$\eps$, pairing gives rate $\tfrac12$, and for the relations of real machines
in general no rate above $\tfrac12$ is achieved and no useful bound below one is proved. Conditions on head motion that
ignore the data are met by a machine that transposes nothing. Bounds for codes that ignore the machine fail, because
abstract relations with the two properties used so far, bounded degree and paired-distance separation, carry codes of
rate tending to one. Counting passes or cuts saturates one scale at a time. The common reason is that the $\log d$ of
the folklore algorithm is a sum over address scales, and every tool that applies to all machines looks at one scale or
at none; the logarithm is proved only where either the coding or the geometry is tied down.

The theorems concern the swap $\sw$ on $V\otimes V$, $V=F^d$, and the code $\ker(\sw-K)$ of matrices that an operator
$K=\sum_aA_a\otimes B_a$ transposes correctly. \emph{Order}: for every complete flag of $V$, the code has dimension at most
the sum of the dimensions of the images of the flag under the $B_a$, whatever the $A_a$ are; so two terms, commuting
families, and, in characteristic $p$ with $d=p^m$, all affine index maps stay at the dimension $d(d+1)/2$ of the
symmetric matrices. \emph{Symmetry}: if each row of a matrix is a fixed linear function of the matrix seen from that
row's position, the same function on each of $M$ intervals of rows, then the matrix ranges over a space of dimension at
most $(d^2+Md)/2$, on a cycle and on a line, with equality for equal intervals, where the bound is the dimension of the
coarse $M\times M$ block transposition. The proof intersects the code with its own translate. It works in every
characteristic and uses a property of machines that cone relations do not record: translated outputs share their
decoder. Variants price each change of the decoder of a single output entry, from one row to the next, at half a
symbol, where a change of class costs half a row, also in the presence of advice over a finite field; bound decoders
with few modulation frequencies; and show that every address digit must be treated asymmetrically, which one unshuffle
does for all digits at once. For classes of rows that no single shift respects the bound
can fail: at $d=6$ two such classes carry a code of dimension $25>24$, and, by a rule taken from a memo to this paper,
two classes carry codes that exceed the bound by $d/4-1$ at every $d=4k\ge8$, over every field; for every fixed number
$M\ge2$ of classes the gain over the bound is unbounded. A later paper of the series, cited in this revision, bounds
two classes of any shape by $\tfrac{41}{64}d^2+O(d^{3/2})$ and shows that the second theorem needs the code, but not
the decoders, to be linear; for sets of matrices and arbitrary decoders two bounds of this paper remain, the
elimination step of the first theorem and the bound for rules that copy one entry, as instances of the
acyclic-subgraph bound. The same bound shows that a set of matrices over $q$ symbols in which every entry is a
function of at most two entries, at positions that move with the row, has at most $q^{d(d+1)/2+d-1}$ elements, and
that reading any number of entries at boundedly many offsets gives no more than $q^{(1/2+o(1))d^2}$; whether three
entries at arbitrary positions of this kind can give more is open.

Both theorems are statements about a linear one-hop model, not about machines, and the question remains open. The
coordinate form of the first theorem, its attainment, the pigeonhole lemma, the second theorem on a cycle for arbitrary
shift sets and on a line, its attainment for equal intervals, its form for changes at single entries with and
without advice, and the twist-class variant are formalized in Lean~4 on the standard axioms. The corollaries that need an algebraically closed
field or the structure of $p$-groups, the copy-rule bound, the modulation and digit corollaries, the passage from
flags to bases, one count in a sharper form of the bound for single entries, the gains for a fixed number of
classes and the bounds for sets are proved in the text and not formalized. Exact computations, at sizes up to six and for the two-class rule and its blow-ups at larger sizes, are
reported as numerical remarks.
